\documentclass[11pt,a4paper]{article}
% preamble.tex — estilo común de todos los apuntes Froth.
% Cada apunte hace % preamble.tex — estilo común de todos los apuntes Froth.
% Cada apunte hace % preamble.tex — estilo común de todos los apuntes Froth.
% Cada apunte hace \input{preamble} justo después de \documentclass.
\usepackage[margin=2.4cm]{geometry}
\usepackage{xcolor}
\definecolor{litblue}{RGB}{31,79,121}
\definecolor{litgreen}{RGB}{27,94,32}
\definecolor{litorange}{RGB}{230,81,0}
\usepackage{parskip}
\usepackage{titlesec}
\titleformat{\section}{\Large\bfseries\color{litblue}}{\thesection}{0.6em}{}
\titleformat{\subsection}{\large\bfseries\color{litblue}}{\thesubsection}{0.6em}{}
\usepackage{tcolorbox}
\usepackage{fancyhdr}
\pagestyle{fancy}
\fancyhf{}
\fancyhead[L]{\small\color{litblue}Froth — Apuntes de ML}
\fancyhead[R]{\small\thepage}
\renewcommand{\headrulewidth}{0.4pt}
\usepackage[colorlinks=true,linkcolor=litblue,urlcolor=litblue]{hyperref}

% Cabecera de cada apunte: \cabecera{numero}{titulo}
\newcommand{\cabecera}[2]{%
  \begin{tcolorbox}[colback=litblue,coltext=white,colframe=litblue,arc=2mm]
    {\Large\bfseries Apunte #1 — #2}\\[3pt]
    {\small Proyecto Froth · aprender Machine Learning construyendo}
  \end{tcolorbox}\vspace{0.8em}}

% Cajas temáticas
\newtcolorbox{concepto}[1]{colback=blue!4,colframe=litblue,title={Concepto: #1},arc=1.5mm}
\newtcolorbox{practica}{colback=green!5,colframe=litgreen,title={Pruébalo tú (teclado en mano)},arc=1.5mm}
\newtcolorbox{ojo}{colback=orange!8,colframe=litorange,title={Ojo / error típico},arc=1.5mm}
\newtcolorbox{resumen}{colback=gray!8,colframe=gray!60,title={En una frase},arc=1.5mm}

\newcommand{\codigo}[1]{\texttt{#1}}
 justo después de \documentclass.
\usepackage[margin=2.4cm]{geometry}
\usepackage{xcolor}
\definecolor{litblue}{RGB}{31,79,121}
\definecolor{litgreen}{RGB}{27,94,32}
\definecolor{litorange}{RGB}{230,81,0}
\usepackage{parskip}
\usepackage{titlesec}
\titleformat{\section}{\Large\bfseries\color{litblue}}{\thesection}{0.6em}{}
\titleformat{\subsection}{\large\bfseries\color{litblue}}{\thesubsection}{0.6em}{}
\usepackage{tcolorbox}
\usepackage{fancyhdr}
\pagestyle{fancy}
\fancyhf{}
\fancyhead[L]{\small\color{litblue}Froth — Apuntes de ML}
\fancyhead[R]{\small\thepage}
\renewcommand{\headrulewidth}{0.4pt}
\usepackage[colorlinks=true,linkcolor=litblue,urlcolor=litblue]{hyperref}

% Cabecera de cada apunte: \cabecera{numero}{titulo}
\newcommand{\cabecera}[2]{%
  \begin{tcolorbox}[colback=litblue,coltext=white,colframe=litblue,arc=2mm]
    {\Large\bfseries Apunte #1 — #2}\\[3pt]
    {\small Proyecto Froth · aprender Machine Learning construyendo}
  \end{tcolorbox}\vspace{0.8em}}

% Cajas temáticas
\newtcolorbox{concepto}[1]{colback=blue!4,colframe=litblue,title={Concepto: #1},arc=1.5mm}
\newtcolorbox{practica}{colback=green!5,colframe=litgreen,title={Pruébalo tú (teclado en mano)},arc=1.5mm}
\newtcolorbox{ojo}{colback=orange!8,colframe=litorange,title={Ojo / error típico},arc=1.5mm}
\newtcolorbox{resumen}{colback=gray!8,colframe=gray!60,title={En una frase},arc=1.5mm}

\newcommand{\codigo}[1]{\texttt{#1}}
 justo después de \documentclass.
\usepackage[margin=2.4cm]{geometry}
\usepackage{xcolor}
\definecolor{litblue}{RGB}{31,79,121}
\definecolor{litgreen}{RGB}{27,94,32}
\definecolor{litorange}{RGB}{230,81,0}
\usepackage{parskip}
\usepackage{titlesec}
\titleformat{\section}{\Large\bfseries\color{litblue}}{\thesection}{0.6em}{}
\titleformat{\subsection}{\large\bfseries\color{litblue}}{\thesubsection}{0.6em}{}
\usepackage{tcolorbox}
\usepackage{fancyhdr}
\pagestyle{fancy}
\fancyhf{}
\fancyhead[L]{\small\color{litblue}Froth — Apuntes de ML}
\fancyhead[R]{\small\thepage}
\renewcommand{\headrulewidth}{0.4pt}
\usepackage[colorlinks=true,linkcolor=litblue,urlcolor=litblue]{hyperref}

% Cabecera de cada apunte: \cabecera{numero}{titulo}
\newcommand{\cabecera}[2]{%
  \begin{tcolorbox}[colback=litblue,coltext=white,colframe=litblue,arc=2mm]
    {\Large\bfseries Apunte #1 — #2}\\[3pt]
    {\small Proyecto Froth · aprender Machine Learning construyendo}
  \end{tcolorbox}\vspace{0.8em}}

% Cajas temáticas
\newtcolorbox{concepto}[1]{colback=blue!4,colframe=litblue,title={Concepto: #1},arc=1.5mm}
\newtcolorbox{practica}{colback=green!5,colframe=litgreen,title={Pruébalo tú (teclado en mano)},arc=1.5mm}
\newtcolorbox{ojo}{colback=orange!8,colframe=litorange,title={Ojo / error típico},arc=1.5mm}
\newtcolorbox{resumen}{colback=gray!8,colframe=gray!60,title={En una frase},arc=1.5mm}

\newcommand{\codigo}[1]{\texttt{#1}}

\begin{document}
\cabecera{06}{Cargar SPECTER y tu primer vector (usar vs.\ construir un modelo)}

\section{Qué hicimos}

En el Apunte 05 vimos la \emph{idea} del embedding. Aquí lo tocamos de verdad: tomamos un
paper real de tu tabla y lo convertimos en su vector.

\begin{enumerate}
  \item Cargamos un paper (título + abstract) desde \codigo{2\_Datos/processed/papers.parquet}.
  \item Cargamos el modelo con una línea: \codigo{SentenceTransformer("allenai-specter")}.
  \item Lo convertimos: \codigo{model.encode(texto)} $\rightarrow$ un vector de \textbf{768 números}.
\end{enumerate}

El paper de prueba fue \emph{``Electrostatically Controlled Enrichment of Lepidolite via
Flotation''} y su vector empezaba por \codigo{[-0.59, -0.16, 1.34, ...]}. Ese paper es ahora
un punto en el mapa del significado.

\begin{concepto}{modelo preentrenado, en la práctica}
La primera vez que llamas a \codigo{SentenceTransformer("allenai-specter")}, la librería
\textbf{descarga} el modelo ($\sim$440 MB) desde el repositorio Hugging Face y lo guarda en
una \emph{caché} en tu disco. Las siguientes veces lo carga de ahí, al instante, sin internet.
No entrenas nada: usas pesos que otros ya calcularon.
\end{concepto}

\begin{ojo}
\textbf{Título + abstract, no solo abstract.} SPECTER fue entrenado para recibir el título
pegado al abstract, porque el título aporta señal fuerte del tema. Por eso el código hace
\codigo{texto = titulo + ". " + abstract} antes de vectorizar. Un detalle pequeño que mejora
los resultados.
\end{ojo}

\begin{ojo}
\textbf{Los \emph{warnings} que salieron son inofensivos.} Uno hablaba de ``symlinks'' (una
optimización de disco en Windows) y otro de ``HF\_TOKEN'' (una llave opcional para descargar
más rápido). Ninguno impide nada: el modelo cargó y funcionó. Aprende a distinguir un
\emph{warning} (aviso, sigue adelante) de un \emph{error} (se detiene).
\end{ojo}

\section{La pregunta clave: ¿construimos nosotros el modelo?}

\textbf{No.} Y conviene tener clarísima la diferencia entre dos cosas que suenan parecidas:

\begin{center}
\begin{tabular}{p{4.3cm} p{4.3cm} p{4.3cm}}
 & \textbf{USAR un modelo} & \textbf{ENTRENAR un modelo} \\[2pt]
\hline
¿Qué haces? & Descargas uno ya hecho y le pasas texto & Le enseñas desde cero con datos; ajustas sus ``pesos'' \\[6pt]
¿Cuánto cuesta? & Una línea de código, segundos & Semanas, muchos datos, tarjetas gráficas (GPU) \\[6pt]
¿En Froth? & \textbf{Fases 2 a 5} (lo de ahora) & \textbf{Fase 6}, opcional y avanzada \\
\end{tabular}
\end{center}

\begin{concepto}{por qué reutilizar es lo correcto (ahora)}
Para dibujar las burbujas y el mapa, el SPECTER genérico —entrenado con millones de papers
de todas las áreas— es más que suficiente. Entrenar tu propio modelo de embeddings desde
cero sería reinventar la rueda, carísimo y sin ganancia clara en esta etapa. Regla del
proyecto: \emph{reciclar sin vergüenza}. Lo original de tu trabajo está en \textbf{cómo
combinas} las piezas y el problema que resuelves, no en re-fabricar el motor.
\end{concepto}

\section{Entonces, ¿cuándo sí tocamos ``el modelo''?}

En la \textbf{Fase 6} (la última, opcional). Ahí el experimento con potencial de tesis es
\emph{afinar} (fine-tuning) un modelo con papers de tu propio dominio (minerales, flotación,
metalurgia) y medir si produce \textbf{mejores} clusters que el SPECTER genérico. Eso sí es
``construir algo tuyo'', y para entonces habrás aprendido PyTorch y cómo funciona un modelo
por dentro. Es el postre, no el plato principal.

\begin{resumen}
El modelo de vectores (SPECTER) ya está construido: solo lo descargamos y usamos, no lo
entrenamos. Usar $\neq$ entrenar; entrenar lo tuyo es la Fase 6, opcional. Por ahora,
tranquilidad: el motor ya está hecho.
\end{resumen}

\end{document}
